%%%%%%%%%%%%%%%%%%%%%%%%%%%%%%%%% LAB-5 %%%%%%%%%%%%%%%%%%%%%%%%%%%%%%%%%%
%>>>>>>>>>>>>>>>>>>>>>>>>>> ПЕРЕМЕННЫЕ >>>>>>>>>>>>>>>>>>>>>>>>>>>>>>>>>>>
%>>>>> Информация о кафедре
%\newcommand{\year}{2021 г.}  % Год устанавливается автоматически
\newcommand{\city}{Санкт-Петербург}  %  Футер, нижний колонтитул на титульном листе
\newcommand{\university}{Национальный исследовательский университет ИТМО}  % первая строка
\newcommand{\department}{Факультет программной инженерии и компьютерной техники}  % Вторая строка
\newcommand{\major}{Направление системного и прикладного программного обеспечения}  % Треьтя строка
%<<<<< Информация о кафедре

%>>>>> Назание работы
\newcommand{\lab}{Лабораторная работа}
\newcommand{\labnumber}{№ 2}                    % порядковый номер работы
\newcommand{\subject}{Основы профессиональной деятельности}         % учебный предмет
\newcommand{\labtheme}{Исследование работы БЭВМ}            % Тема лабораторной работы
\newcommand{\variant}{№ 1033}                % номер варианта работы

\newcommand{\student}{Тюрин Иван Николаевич}    % определение ФИО студента
\newcommand{\studygroup}{P3110}                 % определение учебной группы 
\newcommand{\teacher}{Клименков С. В.,\\[1mm]     % ФИО лектора
                        Ларочкин Г. И.}          % ФИО практика
%<<<<<<<<<<<<<<<<<<<<<<<<<< ПЕРЕМЕННЫЕ <<<<<<<<<<<<<<<<<<<<<<<<<<<<<<<<<<<


%>>>>>>>>>>>>>>>>>>>>>> ПРЕАМБУЛА >>>>>>>>>>>>>>>>>>>>>>>>>

%>>>>>>>>>>>>>>>>>> ПРЕАМБУЛА >>>>>>>>>>>>>>>>>>>>
\documentclass[14pt,final,oneside]{extreport}% класс документа, характеристики
%>>>>> Разметка документа
\usepackage[a4paper, mag=1000, left=3cm, right=1.5cm, top=1cm, bottom=2cm, headsep=0.7cm, footskip=1cm]{geometry} % По ГОСТу: left>=3cm, right=1cm, top=2cm, bottom=2cm,
\linespread{1} % межстройчный интервал по ГОСТу := 1.5
%<<<<< Разметка документа

%>>>>> babel c языковым пакетом НЕ должны быть первым импортируемым пакетом
\usepackage[utf8]{inputenc}
\usepackage[T1,T2A]{fontenc}
\usepackage[russian]{babel}
%<<<<<

%\usepackage{cmap} %поиск в pdf

%>>>...>> прочие полезные пакеты
\usepackage{amsmath,amsthm,amssymb}
\usepackage{mathtext}
\usepackage{indentfirst}
\usepackage{graphicx}
\usepackage{float}
\graphicspath{{/home/ivan/itmo/informatics/latex}}
\DeclareGraphicsExtensions{.pdf,.png,.jpg}
%\usepackage{bookmark}

\usepackage[dvipsnames]{xcolor}
\usepackage{hyperref}  % Использование ссылок
\hypersetup{%  % Настройка разметки ссылок
    colorlinks=true,
    linkcolor=blue,
    filecolor=magenta,      
    urlcolor=magenta,
    %pdftitle={Overleaf Example},
    %pdfpagemode=FullScreen,
}
\usepackage{longtable}
\usepackage{diagbox}
\usepackage[letterspace=150]{microtype} % Спэйсинг (межбуквенный интервал для саголовка) \lsstyle

%>>> верстка в 2 колонки
\usepackage{multicol} % многоколоночная верстка
\setlength{\columnsep}{.15\textwidth} % определение ширины разделителя между колонками

%> кастомный разделитель колонок
\usepackage{tikz} % пакет для векторной графики, чтобы рисовать красивый разделитель колонок
\usetikzlibrary{arrows.meta,decorations.pathmorphing,backgrounds,positioning,fit,petri}
\usepackage{multicolrule} % Для кастомизации разделителя колонок
\SetMCRule{                     % кастомизация разделителя колонок multicolrule
    width=2pt,
    custom-line={               % Tikz код для кастомизации линии разделителя
        \draw [                 % Рисовать
            decorate,           % декорированную (требуются спец настройки пакетов tikz (см. импорт выше)
            decoration={        % вид декорирования
                snake, % Тип - змейка (волнистая)
                amplitude=.5mm, % ширина волн
                pre length=0mm, % участок прямой линии от начала
                %segment length=0mm, % учасок волнистой линии
                post length=0mm % участок прямой линии от конца
            },
            line width=1pt,
            step=10pt
        ] 
        (TOP) to (BOT); % сверху и до низа колонки
    }, 
    extend-top=-5pt, % Вылезти за верхнюю границу колонки 
    extend-bot=-7pt % Вылезти за нижнюю границу колонки  
}
%< кастомный разделитель колонок
%<<< верстка в 2 колонки

%>>>>> Использование листингов
\usepackage{listings} 
\usepackage{caption}
\DeclareCaptionFont{white}{\color{white}} 
\DeclareCaptionFormat{listing}{\colorbox{gray}{\parbox{\textwidth}{#1#2#3}}}

\captionsetup[lstlisting]{format=listing,labelfont=white,textfont=white} % Настройка вида описаний
\lstset{  % Настройки вида листинга
inputencoding=utf8, extendedchars=\true, keepspaces = true, % поддержка кириллицы и пробелов в комментариях
language={},            % выбор языка для подсветки (здесь это Pascal)
basicstyle=\small\sffamily, % размер и начертание шрифта для подсветки кода
numbers=left,               % где поставить нумерацию строк (слева\справа)
numberstyle=\tiny,          % размер шрифта для номеров строк
stepnumber=1,               % размер шага между двумя номерами строк
numbersep=5pt,              % как далеко отстоят номера строк от подсвечиваемого кода
backgroundcolor=\color{white}, % цвет фона подсветки - используем \usepackage{color}
showspaces=false,           % показывать или нет пробелы специальными отступами
showstringspaces=false,     % показывать илигнет пробелы в строках
showtabs=false,             % показывать или нет табуляцию в строках
frame=single,               % рисовать рамку вокруг кода
tabsize=2,                  % размер табуляции по умолчанию равен 2 пробелам
captionpos=t,               % позиция заголовка вверху [t] или внизу [b] 
breaklines=true,            % автоматически переносить строки (да\нет)
breakatwhitespace=false,    % переносить строки только если есть пробел
escapeinside={\%*}{*)}      % если нужно добавить комментарии в коде
}
%<<<<< Использование листингов
\usepackage{csvsimple} %импорт содержимого таблицы из csv


\sloppy % Решение проблем с переносами (с. 119 книга Львовского)
\emergencystretch=25pt


%>>>>>>>>>>>>>>>> ДОПОЛНИТЕЛЬНЫЕ КОМАНДЫ {Для соответствия ГОСТ} >>>>>>>>>>>>>>
%>>>>>> математические функции для удобства
\newcommand{\eps}{\varepsilon}
\newcommand{\limit}{\displaystyle\lim}
\newcommand{\oo}{\infty}
\newcommand{\Dx}{\Delta x}
\newcommand{\cd}{\cdot}
% \newcommand{\note}[2]{\overbrace{#1}^{#2}} % скобка сверху для комментария

%>>>>> Скобки
\newcommand{\lt}{\left}
\newcommand{\rt}{\right}
%<<<<< Скобки

%>>>>> Дроби
\newcommand{\cf}{\cfrac}
\newcommand{\fr}{\frac}
%<<<<< Дроби


%>>>>> Стрелки
\newcommand{\Rarr}{\Rightarrow}% ⇒ следствие
\newcommand{\LRarr}{\Leftrightarrow}% равносильно
\newcommand{\rarr}{\xrightarrow{}}% → стрелка вправо
\newcommand{\nwarr}{\nwarrow}% ↖ север-запад стрелка
\newcommand{\nearr}{\nearrow}% ↗ север-восток стрелка
\newcommand{\swarr}{\swarrow}% ↙ юг-запад стрелка
\newcommand{\searr}{\searrow}% ↘ юг-восток стрелка

\newcommand{\raises}{\nwarrow}% возрастает
\newcommand{\falls}{\swarrow}% убывает
\newcommand{\increases}{\nwarrow}% возрастает
\newcommand{\decreases}{\swarrow}% убывает

%{{{
\makeatletter
\newcommand{\xLeftrightarrow}[2][]{\ext@arrow 0359\Leftrightarrowfill@{#1}{#2}}% следствие с надписью
\makeatother
%}}}

%{{{
\makeatletter
\newcommand{\xRightarrow}[2][]{\ext@arrow 0359\Rightarrowfill@{#1}{#2}}% равносильность с надписью
\makeatother
%}}}
%<<<<< Стрелки
%>>>>> Стиль текста
\newcommand{\hex}[1]{\texttt{0{\footnotesize{x}}#1}}
\newcommand{\ttt}[1]{\texttt{#1}}
%<<<<< Стиль текста

%<<<<<< математические функции для удобства

\newcommand\Chapter[3]{%
    % Принимает 3 аргумента - название главы и дополнительный заголовок и ширина загловка (можно ничего)
    \refstepcounter{chapter}%
    \chapter*{%
        %\hfill % заполнение отступом пространства до заголовка
        \begin{minipage}{#3\textwidth} % Можно изменить ширину министраницы (заголовка)
            \flushleft % Выранивание заголовка по левому краю параграфа (заголовка)
            %\flushright % Выранивание заголовка по правому краю параграфа (заголовка)
            \begin{huge}%
                % Отключена нумерация глав в тексте:
                %:=% \textbf{\chaptername\ \arabic{chapter}\\}
                \textbf{#1}% Первый заголовок
            \end{huge}%
            \\% Перенос сторки
            \begin{Huge}
                #2% Второй заголовок
            \end{Huge}
        \end{minipage}
    }%
    % Отключена нумерация для chapter в toc (table of contents), т.е. Оглавлении (Содержании):
    %:=% \addcontentsline{toc}{chapter}{\arabic{chapter}. #1}
    % Представление главы в содержании:
    \addcontentsline{toc}{chapter}{#1. #2}%
}

\newcommand\Section[1]{
    % Принимает 1 аргумент - название секции
    \refstepcounter{section}
    \section*{%
        \raggedright
        % Отключена дополнительная нумерация chapter в section в тексте документа:
        %:=% \arabic{chapter}.\arabic{section}. #1}
        % Отключена любая нумарация section в тексте документа: (убрать \arabic{section}, оставить название секции)
        \arabic{section}. #1
    }
    
    % Отключена дополнительная нумерация chapter в section в toc (table of contents) Оглавлении (Содержании):
    %:=% \addcontentsline{toc}{section}{\arabic{chapter}.\arabic{section}. #1}
    \addcontentsline{toc}{section}{\arabic{section}. #1} 
}


\newcommand\Subsection[1]{
    % Принимает 1 аргумент - название подсекции
    \refstepcounter{subsection}
    \subsection*{%
        \raggedright%
        % Отключена дополнительная нумерация chapter в section в тексте документа (можно добавить отступ с помощью \hspace*{12pt}):
        %:=% \arabic{chapter}.\arabic{section}.\arabic{subsection}. #1}
        \arabic{section}. \arabic{subsection}. #1
    }
    % Отключена дополнительная нумерация chapter в section в Оглавлении (Содержании):
    %\addcontentsline{toc}{subsection}{\arabic{chapter}.\arabic{section}.\arabic{subsection}. #1}
    \addcontentsline{toc}{subsection}{\arabic{subsection}. #1}
}


\newcommand\Figure[4]{
    % Принимает 4 аргумента - название файла изображения, ее размер в тексте, описание, лэйбл (псевдоним в формате "fig:name") 
    %
    \refstepcounter{figure}
    \begin{figure}[H] %- \usepackage {float} %[h]
        \begin{center}
            \includegraphics[width=#2]{#1}
        \end{center}
        \caption{#3}
        \label{fig:#4}
    \end{figure}
}


\newcommand\Table[3]{
    % Принимает 3 аргумента --- лэйбл name(#1) (псевдоним в формате "tab:name"), ее описание(#2), содержание таблицы(#3) 
    %
    \refstepcounter{table}
    \renewcommand{\arraystretch}{1} % Установка высоты строки таблицы по умолчанию (=1), увеличенное на 0.2 пункта
    % \refstepcounter{table}% увеличение счетчика таблиц
    \begin{table}[htpb]% "right here"=H, "top", "new page", "bottom"
        \label{tab:#1}% лэйбл таблицы, для ссылок
        \resizebox{\columnwidth}{!}{% сжимает очень широкие таблицы, чтобы вместить на страницу
             #3% Содержание таблицы
        }
        % 
        \caption{#2}% Описание стандартными средствами для используемого окружения (table)
        % \captionof{table}{#2}% Описание стандартными средствами
        % \captionof*{figure}{\flushleft \textsc\textbf{Рис. 1.}}% Описание стандартными средствами, как рисунка
        %
        %%> кастомное описание
        % \begin{flushleft}% Кастомное описание
        %     % \textsf{%
        %         \textbf{%
        %             \\[2mm]
        %             #2% Описание к картинке
        %         }%
        %         % \\[8mm]% Отступ
        %     % }%
        % \end{flushleft}
        %%< кастомное описание
    \end{table}
    \renewcommand{\arraystretch}{1} % возврат установка высоты строки таблицы по умолчанию на 1
}


\newcommand\CustomFigure[4]{ % multicols не умеют в table и figure, поэтому приходится извращаться % вставка таблицы с меткой рисунка
    % Принимает 4 аргумента - название файла изображения, ее размер в тексте, описание, лэйбл (псевдоним в формате "fig:name") 
    %
    \refstepcounter{figure}
    \begin{figure}[ht]% "here", "top"
        \begin{center}
            \includegraphics[width=#2]{#1}
        \end{center}
        %
        %\caption{#3}
        \captionof{figure}{#3}% описание стандартными средствами
        % \begin{center}
        \begin{flushleft} % Кастомное описание
            \textbf{%
                #3% Текст описания
            }
        \end{flushleft}
        % \end{center}
        %
        \label{fig:#4}% Лэйбл, для ссылок
    \end{figure}
}


\newcommand\CustomTableFigure[3]{ % multicols не умеют в table и figure, поэтому приходится извращаться % вставка таблицы с меткой рисунка
    %
    % Принимает 3 аргумента --- лэйбл name(#1) (псевдоним в формате "tab:name"), ее описание(#2), содержание таблицы(#3) 
    %
    \begin{center}
        \refstepcounter{figure}
        \label{tab:#1}% лэйбл таблицы, для ссылок
        \resizebox{\columnwidth}{!}{% сжимает очень широкие таблицы, чтобы вместить на страницу
            #3% Содержание таблицы
        }
        % 
        \captionof{figure}{#2}% Описание стандартными средствами
        % \captionof*{figure}{\flushleft \textsc\textbf{Рис. 1.}}% Описание стандартными средствами
        %
        \begin{flushleft}% Кастомное описание
            % \textsf{%
                \textbf{%
                    \\[2mm]
                    #2% Описание к картинке
                }%
                % \\[8mm]% Отступ
            % }%
        \end{flushleft}
    \end{center}
}


\newcommand{\InkscapeFigure}[4]{% Вставки иллюстраций из Inkscape (pdf+latex)
    %
    % Принимает 4 параметра: #1 название файла, #2 описание, #3 лейбл #4 размер
    %
    % \begin{minipage}{#4}
        \begin{figure}[htbp]
            \centering
            \def\svgwidth{#4}
            \import{./figures/}{#1.pdf_tex}
            \caption{#2}
            \label{fig:#3}
        \end{figure}
    % \end{minipage}
}


\newcommand\Equation[3]{% Кастомное оформление выражений
    %
    % Принимает 3 аргумента --- лэйбл name (#1) (псевдоним в формате "tab:name"), его описание(#2), содержание выражения (#3) 
    %
    \textbf{#2}% описание
    \begin{equation}
        #3% содержимое выражений
        \label{eq:#1}% лэйбл
    \end{equation}
}

%<<<<<<<<<<<<<<<<<<<<<<<<<<<< ДОПОЛНИТЕЛЬНЫЕ КОМАНДЫ <<<<<<<<<<<<<<<<<<<<<<<<<<


\newcommand{\Rarr}{\Rightarrow}
\newcommand{\LRarr}{\Leftrightarrow}
\newcommand{\rarr}{\xrightarrow{}}
\newcommand{\rt}{\right}
\newcommand{\lt}{\left}

\makeatletter
\newcommand{\xLeftrightarrow}[2][]{\ext@arrow 0359\Leftrightarrowfill@{#1}{#2}}
\makeatother

\makeatletter
\newcommand{\xRightarrow}[2][]{\ext@arrow 0359\Rightarrowfill@{#1}{#2}}
\makeatother
%<<<<<<<<<<<<<<<<<<<<<< ПРЕАМБУЛА <<<<<<<<<<<<<<<<<<<<<<<<<



%%%%%%%%%%%%%%%%%%% СОДЕРЖИМОЕ ОТЧЕТА %%%%%%%%%%%%%%%%%%%%%
%>>>>>>>>>>>>>>> ''''''''''''''''''''''' >>>>>>>>>>>>>>>>>>
\begin{document}


%>>>>>>>>>>>>>>>> ОПРЕДЕЛЕНИЕ НАЗВАНИЙ >>>>>>>>>>>>>>>>>>>>
% Переоформление некоторых стандартных названий
%\renewcommand{\chaptername}{Лабораторная работа}
\renewcommand{\chaptername}{\lab\ \labnumber} % переименование глав
\def\contentsname{Содержание} % переименование оглавления
%<<<<<<<<<<<<<<<< ОПРЕДЕЛЕНИЕ НАЗВАНИЙ <<<<<<<<<<<<<<<<<<<<


%>>>>>>>>>>>>>>>>> ТИТУЛЬНАЯ СТРАНИЦА >>>>>>>>>>>>>>>>>>>>>
%>>>>>>>>>>>>>>>>>>> ТИТУЛЬНЫЙ ЛИСТ >>>>>>>>>>>>>>>>>>>>>>>
\begin{titlepage}

    % Название университета
    \begin{center}
    \textsc{%
        \university\\[5mm]
        \department\\[2mm]
        \major\\
    }

    \vfill
    \vfill
    % Название работы
    \textbf{\reporttype\ \labnumber\\[3mm]
    курса <<\subject>> \\[6mm]
    по теме: <<\labtheme>>\\[3mm]
    Вариант \variant\\[20mm]
    }
    \end{center}


\hfill
% Информация об авторе работы и проверяющем
\begin{minipage}{.5\textwidth}
    \begin{flushright}
        
            
        Выполнил студент:\\[2mm] 
        \student\\[2mm]
        группа: \studygroup\\[5mm]

        Преподаватель:\\[2mm] 
        \teacher

    \end{flushright}
\end{minipage}

\vfill

    % Нижний колонтитул первой страницы
    \begin{center}
        \city, \the\year\,г.
    \end{center}

\end{titlepage}
%<<<<<<<<<<<<<<<<<<< ТИТУЛЬНЫЙ ЛИСТ <<<<<<<<<<<<<<<<<<<<<<<


%<<<<<<<<<<<<<<<<< ТИТУЛЬНАЯ СТРАНИЦА <<<<<<<<<<<<<<<<<<<<<


%>>>>>>>>>>>>>>>>>>>>> СОДЕРЖАНИЕ >>>>>>>>>>>>>>>>>>>>>>>>>
% Содержание
\tableofcontents
%<<<<<<<<<<<<<<<<<<<<< СОДЕРЖАНИЕ <<<<<<<<<<<<<<<<<<<<<<<<<


%%%%%%%%%%%%%%%%%%%%%%% КОД РАБОТЫ %%%%%%%%%%%%%%%%%%%%%%%%
%>>>>>>>>>>>>>>>>>>>'''''''''''''''''>>>>>>>>>>>>>>>>>>>>>
\newpage
\Chapter{\lab\ \labnumber}{\labtheme}{}
\Section{Задание варианта \variant}
\begin{center}
    , , ,
\end{center}

\indent
По выданному преподавателем варианту определить функцию, вычисляемую программой, область представления и область допустимых значений исходных данных и результата, выполнить трассировку программы, предложить вариант с меньшим числом команд. При выполнении работы представлять результат и все операнды арифметических операций знаковыми числами, а логических операций набором из шестнадцати логических значений.\\

\begin{center}
    \includegraphics[width=0.15\textwidth]{task}
\end{center}

\begin{center}
    ' ' ' 
\end{center}

\Section{Описание программы}

\Subsection{Назначение программы и реализуемые ею функция (формула)}

Описание программы представлено в таблице \ref{tab:program}.\\
Смысловую часть программы можно изобразить формулой:\\ $R=A\And(C-B)$. Программа считает разность переменных $C$ и $B$, затем логически умножает полученное значение на $A$.

\begin{table}[htbp]
    \renewcommand{\arraystretch}{1.2} % Установка высоты строки таблицы
    \centering
    \resizebox{\columnwidth}{!}{%
    \begin{tabular}{|l|l|l|l|}
            \hline
           Адрес, (0x)   & Данные/      & Мнемоника & Результат\\
                         & /Команда, (0x)&           &           \\\hline
           075 &    2075 & --           & Данные (A)\\\hline
           076 &    A082 & --           & Данные (B)\\\hline
           077 &    A082 & LD 0x082     & AC = C \\\hline
           078 &    6076 & SUB 0x076    & AC = C -- B \\\hline
           079 &    E080 & ST 0x080     & Сохранение промеж. рез.:\\
               &         &              & D = C -- B\\\hline
           07A &    0200 & CLA          & AC = 0x0000 = 0\\\hline
           07B &    0280 & NOT          & AC = \textasciicircum AC = 0xFFFF = --1 \\\hline
           07C &    2075 & AND 0x075    & AC = 0xFFFF$\And$A = A \\\hline
           07D &    2080 & AND 0x080    & AC = A$\And$D \\\hline
           07E &    E081 & ST 0x081     & Сохранение промеж. рез.:\\
               &         &              & R = A$\And$D \\\hline
           07F &    0100 & HLT          & Отключение ТГ,\\
               &         &              & переход в пульт. реж. \\\hline
           080 &    0100 & --           & Промеж. данные:\\
               &         &              & D = C -- B \\\hline
           081 &    6076 & --           & Результат:\\
               &         &              & R = A$\And$D = A$\And$(C -- B)\\\hline
           082 &    2080 & --           & Данные (C) \\\hline
    \end{tabular}
    }
    \caption{Описание данной программы}
    \label{tab:program}
\end{table}


\Subsection{Расположение в памяти ЭВМ программы, исходных данных и результатов}

Программа в памяти ЭВМ распологается в ячейках памяти с адресами 0x077 -- 0x07F. Исходные данные находятся по адресам 0x075, 0x076, 0x082. Результат программы находится по адресу 0x081. 

\Subsection{Адреса первой и последней выполняемой команд программы}
Адрес первой выполняемой команды -- 0x077. Адрес последней выполняемой команды (HLT) -- 0x7F. 

\Subsection{Область представления исходных данных и результата}
A, R -- набор из 16 однобитовых значений. B, С -- знаковые 16-ти разрядные числа.  Результат (A + C) -- набор из 16 однобитовых значений.

\Subsection{Область допустимых значений исходных данных и результата}
\begin{equation}
\text{Установим свойство: }
\lt[N\leqslant \text{\textasciicircum}X \leqslant M\rt] \LRarr 
\lt[\text{\textasciicircum}M\leqslant X \leqslant \text{\textasciicircum}N\rt]\label{eq:not}
\end{equation}


$$
A\And(C-B)\Rarr
\begin{array}{cc}
 & \boxed{A_i\in \lbrace 1, 0 \rbrace, \text{ где } 0\leqslant i \leqslant 15};\\[2mm]
& \boxed{X=(C-B),\; X_i\in \lbrace 1, 0 \rbrace, \text{ где } 0\leqslant i \leqslant 15};
\end{array}
$$\\[2mm]


$\lt[-2^{15}\leqslant C-B\leqslant2^{15}-1\rt] 
\LRarr
\lt[-2^{15}\leqslant C + (\text{\textasciicircum} B + 1)\leqslant2^{15}-1 \rt] \xRightarrow{\text{Сл. } 1^\circ, 2^\circ, 3^\circ}$\\[2mm]


\textbf{$1^\circ.$} (Смотри рисунок \ref{fig:way1})

$$
\begin{cases}
-2^{14}\leqslant C\leqslant2^{14}-1\\
-2^{14}-1\leqslant \text{\textasciicircum}B\leqslant2^{14}-2\\
\end{cases} 
\xLeftrightarrow{\text{св-во. \eqref{eq:not}}} 
\begin{cases}
-2^{14}\leqslant C\leqslant2^{14}-1\\
-2^{14}+1\leqslant B\leqslant2^{14}-1\\
\end{cases} 
$$\\[5mm]

\textbf{$2^\circ.$} (Смотри рисунок \ref{fig:way2})

$$
\begin{cases}
-2^{15}\leqslant C\leqslant 0\\
-1\leqslant \text{\textasciicircum}B\leqslant2^{15}-2\\
\end{cases} 
\xLeftrightarrow{\text{св-во. \eqref{eq:not}}} 
\begin{cases}
-2^{15}\leqslant C\leqslant 0\\
-2^{15}+1\leqslant B\leqslant -1\\
\end{cases} 
$$\\[5mm]

\textbf{$3^\circ.$} (Смотри рисунок \ref{fig:way3})

$$
\begin{cases}
0 \leqslant C\leqslant 2^{15}-1\\
-2^{15}-1\leqslant \text{\textasciicircum}B\leqslant-1\\
\end{cases} 
\xLeftrightarrow{\text{св-во. \eqref{eq:not}}} 
\begin{cases}
0 \leqslant C\leqslant 2^{15}-1\\
-1 \leqslant B\leqslant 2^{15}-1\\
\end{cases} 
$$\\[5mm]

Таким образом итоговое ОДЗ:\\

$$
\begin{cases}
A_i\in \lbrace 1, 0 \rbrace, \text{ где } 0\leqslant i \leqslant 15\\
\lt[
\begin{array}{l}
\begin{cases}
-2^{14}\leqslant C\leqslant2^{14}-1\\
-2^{14}+1\leqslant B\leqslant2^{14}-1\\
\end{cases}\\[1mm]
\begin{cases}
-2^{15}\leqslant C\leqslant 0\\
-2^{15}+1\leqslant B\leqslant -1\\
\end{cases}\\[1mm]
\begin{cases}
0 \leqslant C\leqslant 2^{15}-1\\
-1 \leqslant B\leqslant 2^{15}-1\\
\end{cases} 
\end{array}
\rt.
\end{cases}
$$

\incfig{lab-2-fopa-ODZ-1}{Случай 1.}{way1}{0.9\textwidth}
\incfig{lab-2-fopa-ODZ-2}{Случай 2.}{way2}{0.9\textwidth}
\incfig{lab-2-fopa-ODZ-3}{Случай 3.}{way3}{0.9\textwidth}

\Section{Таблица трассироваки}
Промежуточные данные после выполнения каждой команды программы представлены в таблице трассировки \ref{tab:trace}.
\begin{table}[htbp]
    \renewcommand{\arraystretch}{1.2} % Установка высоты строки таблицы
    \centering
    \resizebox{\columnwidth}{!}{%
    \begin{tabular}{*{13}{|c}|}
        \hline
        \multicolumn{2}{|c|}{Выполняемая} & \multicolumn{8}{c|}{Содержимое регистров процессора после} & \multicolumn{2}{c|}{Ячейка, содержимое}\\
        \multicolumn{2}{|c|}{команда}     & \multicolumn{8}{c|}{выполнения команды.} & \multicolumn{2}{c|}{которой изменилось после}\\
        \multicolumn{2}{|c|}{}&\multicolumn{8}{c|}{}&\multicolumn{2}{c|}{выполнения команды}\\\hline
        Адрес & Знчение & IP & CR & AR & DR & SP & BR & AC & NZVC & Адрес & Значение\\\hline
        077 & A082 & 077 & 0000 & 000 & 0000 & 000 & 0000  & 004 & 0100 &   & \\
        077 & A082 & 078 & A082 & 082 & 2080 & 000 & 0077  & 000 & 0000 &   & \\
        078 & 6076 & 079 & 6076 & 076 & A082 & 000 & 0078  & 000 & 0000 &   & \\
        079 & E080 & 07A & E080 & 080 & 7FFE & 000 & 0079  & 000 & 0000 & 080 & 7FFE\\
        07A & 0200 & 07B & 0200 & 07A & 0200 & 000 & 007A  & 004 & 0100 &  &\\
        07B & 0280 & 07C & 0280 & 07B & 0280 & 000 & 007B  & 008 & 1000 &  &\\
        07C & 2075 & 07D & 2075 & 075 & 2075 & 000 & 007C  & 000 & 0000 &  &\\
        07D & 2080 & 07E & 2080 & 080 & 7FFE & 000 & 007D  & 000 & 0000 &  &\\
        07E & E081 & 07F & E081 & 081 & 2074 & 000 & 007E  & 000 & 0000 & 081 & 2074\\
        07F & 0100 & 080 & 0100 & 07F & 0100 & 000 & 007F  & 000 & 0000 &  &\\ \hline
    \end{tabular}
    }
    \caption{Таблица трассировки.}
    \label{tab:trace}
\end{table}



\Section{Вариант программы с меньшим числом команд}
Описание программы с меньшим числом команд представлено в таблице \ref{tab:optimized}.
Смысловую часть программы можно изобразить формулой:\\ $R=A\And(C-B)$. Программа считает разность переменных $C$ и $B$, затем логически умножает полученное значение на $A$ и записывает полученное значение в первую свободную ячейку памяти (сразу после команд программы).

\begin{table}[htbp]
    \renewcommand{\arraystretch}{1.2} % Установка высоты строки таблицы
    \centering
    \begin{tabular}{|l|l|l|l|}
            \hline
           Адрес, (0x)   & Данные/      & Мнемоника & Результат\\
                         & /Команда, (0x)&           &          \\\hline
           000 &    2075 & --           & Данные (A)\\\hline
           001 &    A082 & --           & Данные (B)\\\hline
           002 &    2080 & --           & Данные (C)\\\hline
           003 &    A002 & LD 0x002     & AC = C \\\hline
           004 &    6001 & SUB 0x001    & AC = AC - B \\\hline
           005 &    2000 & AND 0x000    & AC = A$\And$AC = A$\And$(C -- B)\\\hline
           006 &    E008 & ST 0x008     & Cохранение рез.: \\
               &         &              & 0x008: A$\And$(C -- B)\\\hline
           007 &    0100 & HLT          & Отключение ТГ,\\
               &         &              & переход в пульт. реж. \\\hline
    \end{tabular}
    \caption{Описание сокращенной программы}
    \label{tab:optimized}
\end{table}





\Section{Вывод}
В ходе выполнения лабораторной работы было изучено устройство БЭВМ, основные её команды. На  эмуляторе БЭВМ была запущена простейшая программа и произведён разбор состояний БЭВМ при  пошаговом выполнении программы.\\
\newpage
%<<<<<<<<<<<<<<<<<<<<<< КОД РАБОТЫ <<<<<<<<<<<<<<<<<<<<<<<<


%>>>>>>>>>>>>>>>> СПИСОК ЛИТЕРАТУРЫ >>>>>>>>>>>>>>>>>>>>>>>
%
\bibliographystyle{plain}

\begin{thebibliography}{3}
    \addcontentsline{toc}{chapter}{Список лиетратуры}

    \bibitem{gutgut:1}
    Код Хэмминга. Пример работы алгоритма. URL: \url{https://habr.com/ru/post/140611/};

    \bibitem{gutgut:2}
    Избыточное кодирование, код Хэмминга. URL: \url{https://neerc.ifmo.ru/wiki/index.php?title=%D0%98%D0%B7%D0%B1%D1%8B%D1%82%D0%BE%D1%87%D0%BD%D0%BE%D0%B5_%D0%BA%D0%BE%D0%B4%D0%B8%D1%80%D0%BE%D0%B2%D0%B0%D0%BD%D0%B8%D0%B5,_%D0%BA%D0%BE%D0%B4_%D0%A5%D1%8D%D0%BC%D0%BC%D0%B8%D0%BD%D0%B3%D0%B0}.

\end{thebibliography}  % Для соответсвия гост, придется доработать. Нужен файл .bib
%<<<<<<<<<<<<<<<<<<<< СПИСОК ЛИТЕРАТУРЫ <<<<<<<<<<<<<<<<<<<


\end{document}
%<<<<<<<<<<<<<<<< ,,,,,,,,,,,,,,,,,,,,,,, <<<<<<<<<<<<<<<<<
%<<<<<<<<<<<<<<<<<<< СОДЕРЖИМОЕ ОТЧЕТА <<<<<<<<<<<<<<<<<<<<
